\documentclass[11pt,a4paper]{article}
\usepackage[margin=22mm]{geometry}
\usepackage{amsmath,amssymb}
\usepackage[hidelinks]{hyperref}
\setlength{\parindent}{0pt}
\setlength{\parskip}{5pt}
\raggedright
\begin{document}
{\Large A merger follows its own queue: local motion and kinetic memory}\\[4pt]
8 October 2026 --- Q.J-A approved; analytical completion, no new law adopted

\textbf{Scope.} Source-free local triangular KW with held desired speeds,
one-sided initial states and an explicitly resolved finite queue encounter.
Records describe restrictions and add \emph{zero} traffic population. Results
are conditional mathematical references, not the complete lane-source/node/E1
solver, a reached metering comparison, clearance proof or gradient evidence.
Source and adoption boundaries are in \texttt{motion\_contact\_source\_audit.md}.

\section*{1. Ordinary traffic already supplies the inactive limit}
With speed units m/s, density in equivalents/m and flow in equivalents/s, let
\begin{equation}
 F(\rho)=\min\{u\rho,w(\kappa-\rho)\},\quad
 \rho_c=\frac{w\kappa}{u+w},\quad
 V(\rho)=\begin{cases}u,&0\le\rho\le\rho_c,\\
 w(\kappa/\rho-1),&\rho_c<\rho\le\kappa.\end{cases}
\end{equation}
Vacuum speed is $V(0)=u$. The ordinary relative flow seen by a path of speed
$v$ is $g_v(\rho)=F(\rho)-v\rho$. Thus
\begin{equation}
 g_u(\rho)\le0\quad(0\le\rho\le\kappa),\qquad
 R(u)=\sup_{0\le\rho\le\kappa}g_u(\rho)=0.
\end{equation}
A zero-passing upper bound along a free-speed path adds no restriction to
ordinary KW. Replacing it by a stationary zero entrance capacity for an entire
step does not preserve this exact limit.

At an initial point take constant left/right traffic states $\rho_L,\rho_R$
and a held desired speed $s\in[0,u]$. Write
\begin{equation}
 c=V(\rho_R),\qquad v=\min\{s,c\}.
\end{equation}
If $c\le s$, the path rides the right stream. For congested $\rho_R$,
a left-to-right shock when $\rho_L<\rho_R$ has speed
\begin{equation}
 \sigma=\frac{F(\rho_R)-F(\rho_L)}{\rho_R-\rho_L}\le c,
 \qquad F(\rho_L)\ge c\rho_L\quad(\rho_L\le\rho_R).
\end{equation}
The latter inequality follows branchwise from $F(\rho)/\rho\ge V(\rho_R)$.
If $\rho_L>\rho_R$ both states are congested and their contact moves at $-w$.
Consequently the ordinary solution's right trace on the path is $\rho_R$,
and $g_c(\rho_R)=0$. The case of vacuum on the left gives the queue rear
itself. For free right traffic with $c=u$, equation (2) covers the degenerate
case. There is no extra service deduction when the path merely follows the
existing stream. An already existing queue is not erased.

\newpage
\section*{2. The active local fan is conservative and one-sided}
If $s<c$, then $s<u$. Zero relative passing and triangular KW give
\begin{equation}
 \rho_U(s)=\frac{w\kappa}{w+s}\ge\rho_c,\quad
 q_U(s)=s\rho_U(s),\quad \rho_D=q_D=0.
\end{equation}
The saturated following trace and vacuum are separated by a path of speed
$s$, since
\begin{equation}
 q_U-s\rho_U=0=q_D-s\rho_D.
\end{equation}
This is the mass jump condition on a moving boundary, not a new population.
Because $s<c$, $\rho_R<\rho_U$. The vacuum-to-right front has speed
\begin{equation}
 \sigma_R=\frac{F(\rho_R)}{\rho_R}=c>s
 \quad(\rho_R>0);
\end{equation}
for $\rho_R=0$ there is simply vacuum. When $\rho_L<\rho_U$, the left shock
is behind the path:
\begin{equation}
 \sigma_L=\frac{q_U-F(\rho_L)}{\rho_U-\rho_L}
 =s-\frac{g_s(\rho_L)}{\rho_U-\rho_L}\le s,
 \qquad g_s(\rho_L)\ge0.
\end{equation}
These density-increasing shocks satisfy the concave-FD entropy condition.
For $\rho_L>\rho_U$, the two congested states join by a contact at $-w$.
For $\rho_L=\rho_U$ no left wave is needed. For $\rho_L=0$, $\sigma_L=s$:
the following saturated state has zero width until traffic catches up. It
does not manufacture a queue or positive service without incoming traffic.
The $s=0$ limit is a stationary jam/void interface, with the same conservation
identity. This gives a local solution with no arbitrary anticipation window.

\section*{3. Same-count coincident records can separate without added mass}
For finitely many records born at the same lane point/count, hold desired
speeds $s_i$ and put $v_i=\min\{s_i,c\}$, $v_0=\min_i v_i$.
If $v_0<c$, the preceding fan for $v_0$ gives vacuum between its path and
the front moving at $c$. All faster records lie in that vacuum or on the
front; their constant-speed segments are compatible with the own field.
If all desired speeds exceed $c$, all records ride the original right stream.
Coincident equal-speed records remain retained.

In this source-free local problem, define $N_x=-\rho$, $N_t=F(\rho)$.
For any two separating paths $\phi_i\le\phi_j$ in the void,
\begin{equation}
 N(t,\phi_i)-N(t,\phi_j)=\int_{\phi_i}^{\phi_j}\rho(t,x)\,dx=0,
 \qquad \frac{d}{dt}N(t,\phi_i)=0.
\end{equation}
No rank spacing, vehicle identifiers or integer vehicle mass is implied.
This establishes compatibility for held local segments; later sources and
kinetic changes require retaining all records, not deleting the faster ones.

\newpage
\section*{4. A restriction catches its own queue}
Take $u=30$, $w=5$, $\kappa=3/20$, an empty gap of length $d=20$ and a
queue ahead with $\rho_J=3/40$. Then
\begin{equation}
 F(\rho_J)=3/8,\qquad c=V(\rho_J)=5,
 \qquad b(t)=20+5t.
\end{equation}
Let an already present restriction start at the entrance with held desired
speed $s=15$. The finite queue extends far enough downstream (e.g. to a free
outlet at 200 m) that its far-end release wave cannot affect the four-second
example. The vacuum/queue rear $b(t)$ is a forward shock. Catch time and
own motion follow from the intersection, not from prescribed external paths:
\begin{equation}
 t_* = \frac{d}{s-c}=2,\qquad
 \phi(t)=\begin{cases}15t,&0\le t\le2,\\
 30+5(t-2),&2\le t\le4.\end{cases}
\end{equation}
Thus
\begin{equation}
 \phi(4)=40,\qquad
 \frac{\phi(4)-\phi(0)}4=10,\qquad v(4^-)=5.
\end{equation}
The four seconds are an analytical illustration window, not a selected CTM
timestep. For any regular interval of length $\Delta t$, an old record with
gap $d=(s-c)\Delta t/2$ reaches the queue halfway through the interval. Such
an old record may lie arbitrarily close to a queued cell face even when
$u\Delta t\le\Delta x$; no CFL violation is needed for an interior contact.

Holding 15 throughout puts the record at 60 m, beyond the queue rear.
Holding 5 puts it at 20 m and brakes too early. The straight path $10t$ has
the correct endpoint but does not reach the queue until second four.

With plentiful ordinary entrance demand behind the record, the following
state before contact is $\rho_U(15)=3/80$, $q_U(15)=9/16$. Upon contact the
following density becomes $\rho_U(5)=3/40=\rho_J$. A congested wave carrying
this change travels backward at $-w$, so
\begin{equation}
 t_{\rm entrance}=t_*+\frac{\phi(t_*)}{w}=2+30/5=8.
\end{equation}
For $t\le4$ the entrance has not yet felt that slowdown. Its cumulative
count and the own population behind the restriction at second four agree:
\begin{equation}
 C(4)=\frac9{16}\,4=\frac94,
 \qquad \int_0^{40}\rho(4,x)\,dx
 =\frac3{80}\,20+\frac3{40}\,20=\frac94.
\end{equation}
The backward wave at second four is at $30-5(4-2)=20$ m. By contrast, a
constant 10 m/s restriction has $\rho_U(10)=1/20$ and $C(4)=2$. The two
paths share an endpoint but differ by $1/4$ equivalent in entrance count.
No control pair or complete-TST difference is inferred. The existing-record,
no-lateral-source diagnostic is not a change to E1 in the full task.

\newpage
\section*{5. The proposed fixed-clock completion needs two different outputs}
The source sampled kinetics use one interval speed in both updates:
\begin{equation}
 \phi^{j+1}=\phi^j+v^j\Delta t,\qquad
 d^{j+1}=\min\{u,v^j+a(v^j,\phi^j)\Delta t\}.
\end{equation}
At an interior contact, an interval-average speed and the ending kinetic
speed are different. An average can encode distance, but cannot also encode
the exact state from which the next acceleration begins. This does not refute
a single-speed \emph{numerical approximation}; it identifies a motion/kinetic
error that such an approximation must qualify. Exact information retention
does not demand an exact continuum field.

\textbf{Candidate K.A (not adopted).} Regular public updates and E1 births
stay at $t_j$. Positive acceleration is refreshed only at those times.
Within the interval let $W_i(t)=V(\rho(t,\phi_i(t)+))$ be the immediate
own ahead trace speed. A braking-only running cap defines, conditionally on
a resolved local field/contact sequence,
\begin{equation}
 v_i(t)=\min\left\{d_i^j,\inf_{t_j\le r\le t}W_i(r)\right\},\qquad
 \phi_i^{j+1}=\phi_i^j+\int_{t_j}^{t_{j+1}}v_i(t)\,dt.
\end{equation}
Instant braking uses the source idealization; the running cap prevents
instantaneous recovery if the queue vanishes later in the same interval.
Keep the true $v_i(t_{j+1}^-)$ and refresh
\begin{equation}
 d_i^{j+1}=\min\{u,v_i(t_{j+1}^-)
       +a_i(v_i(t_{j+1}^-),\phi_i^j)\Delta t\}.
\end{equation}
The old start position in the grade argument preserves the original spatial
phase; the example is level road. Other grade semantics must not be silently
changed. New E1 records are initialized by their approved donor speed, without
this old-record increment or retroactive movement. On a no-contact interval
these equations reproduce (15); on contact intervals they are an explicit
adaptation. They do not by themselves define the accepted-source/node update.

In the example, using average speed 10 instead of final speed 5 would initialize
the next desired speed from the wrong kinetic state. On a level road with
$a(v)=a_0(1-v/v_{\max})$ the difference before the common free-speed cap is
\begin{equation}
 [10+a(10)\Delta t]-[5+a(5)\Delta t]
 =5\left(1-\frac{a_0\Delta t}{v_{\max}}\right).
\end{equation}
This is a kinetic-memory diagnostic, not a claim of a reached control gradient.

\newpage
\section*{6. Continuous kinetics is a distinct alternative}
\textbf{Candidate K.B (not adopted).} Integrate $\dot\phi=v$ and the source
acceleration function continuously while motion is unconstrained. Enforce the
own ahead speed cap by instantaneous braking; after release, positive
acceleration resumes immediately rather than waiting for $t_{j+1}$. Meter
updates and E1 births can remain regular. For level free road put
$\lambda=a_0/v_{\max}$; before reaching $u$,
\begin{equation}
 v(t_j+h)=v_{\max}-(v_{\max}-v(t_j))e^{-\lambda h},\qquad
 \phi(t_j+h)-\phi(t_j)
 =v_{\max}h-\frac{v_{\max}-v(t_j)}{\lambda}(1-e^{-\lambda h}).
\end{equation}
Its free-motion samples differ from the source Euler recurrence even without
a queue. The acceleration family is supported; this integration and the
joined timing contract need explicit adoption. No existing full endogenous
theorem is asserted for either alternative.

\textbf{Candidate K.C (unqualified, not adopted).} Keep the source sampled
single-speed recurrence and the explicit own density-window predictor; use a
moving-relative prescribed-segment field construction rather than the direct
stationary zero mask. Qualify its endogenous motion/wave errors against the
local contact reference. This is an admissible numerical review route, not
record compression. No tolerance, convergence or gradient claim is established.

\section*{7. Where this result sits in the review}
Q.J-A authorized completion work, not adoption of either timing law. The
local activation/motion and coincidence compatibility above are positive
bounded results. The remaining choice K concerns how interior contacts alter
the sampled kinetic state. K.A minimally amends the interval transcription;
K.B also changes unconstrained recovery. Both preserve exact records, M.A,
N.A and E1. No cost or approximation criterion selects between them.

Lateral sources still prevent using a globally constant source-free count:
for $\rho_t+F_x=S$ and $N=C-\int_0^x\rho$, away from point-source ambiguity,
\begin{equation}
 \frac{d}{dt}N(t,\phi(t))
 =F(t,\phi)-\dot\phi\rho(t,\phi)-\int_0^{\phi(t)}S(t,x)\,dx.
\end{equation}
A source on the path needs a one-sided attachment. The completed node must
share receiving once, pair all population transfers and apply E1 afterward.
The current local result supplies none of the missing arbitrary-source or
clearance theorems. Full-population TST, meaningful reached metering leverage,
gradient correctness/usefulness and necessity against aggregate/memory
competitors remain unestablished. Stage 3 stays open; Stage 4 stays unopened.
\end{document}
